\documentclass[10pt,a4paper]{amsart}
\usepackage{amsmath,amssymb,amsthm,graphicx}
\usepackage{verbatim}
\newenvironment{reference}{\comment}{\endcomment}
%\newenvironment{reference}{}{}
\newenvironment{slogan}{\comment}{\endcomment}
\newenvironment{history}{\comment}{\endcomment}

% For commutative diagrams we use Xy-pic
\usepackage[all]{xy}

% We use 2cell for 2-commutative diagrams.
\xyoption{2cell}
% Theorem environments.
%
\theoremstyle{plain}
\newtheorem{theorem}[subsection]{Theorem}
\newtheorem{proposition}[subsection]{Proposition}
\newtheorem{lemma}[subsection]{Lemma}

\theoremstyle{definition}
\newtheorem{definition}[subsection]{Definition}
\newtheorem{example}[subsection]{Example}
\newtheorem{exercise}[subsection]{Exercise}
\newtheorem{situation}[subsection]{Situation}

\theoremstyle{remark}
\newtheorem{remark}[subsection]{Remark}
\newtheorem{remarks}[subsection]{Remarks}

\numberwithin{equation}{subsection}

% Macros
%
\def\lim{\mathop{\mathrm{lim}}\nolimits}
\def\colim{\mathop{\mathrm{colim}}\nolimits}
\def\Spec{\mathop{\mathrm{Spec}}}
\def\Hom{\mathop{\mathrm{Hom}}\nolimits}
\def\Ext{\mathop{\mathrm{Ext}}\nolimits}
\def\SheafHom{\mathop{\mathcal{H}\!\mathit{om}}\nolimits}
\def\SheafExt{\mathop{\mathcal{E}\!\mathit{xt}}\nolimits}
\def\Sch{\mathit{Sch}}
\def\Mor{\mathop{\mathrm{Mor}}\nolimits}
\def\Ob{\mathop{\mathrm{Ob}}\nolimits}
\def\Sh{\mathop{\mathit{Sh}}\nolimits}
\def\NL{\mathop{N\!L}\nolimits}
\def\CH{\mathop{\mathrm{CH}}\nolimits}
\def\proetale{{pro\text{-}\acute{e}tale}}
\def\etale{{\acute{e}tale}}
\def\QCoh{\mathit{QCoh}}
\def\Ker{\mathop{\mathrm{Ker}}}
\def\Im{\mathop{\mathrm{Im}}}
\def\Coker{\mathop{\mathrm{Coker}}}
\def\Coim{\mathop{\mathrm{Coim}}}

% Boxtimes
%
\DeclareMathSymbol{\boxtimes}{\mathbin}{AMSa}{"02}

%
% Macros for moduli stacks/spaces
%
\def\QCohstack{\mathcal{QC}\!\mathit{oh}}
\def\Cohstack{\mathcal{C}\!\mathit{oh}}
\def\Spacesstack{\mathcal{S}\!\mathit{paces}}
\def\Quotfunctor{\mathrm{Quot}}
\def\Hilbfunctor{\mathrm{Hilb}}
\def\Curvesstack{\mathcal{C}\!\mathit{urves}}
\def\Polarizedstack{\mathcal{P}\!\mathit{olarized}}
\def\Complexesstack{\mathcal{C}\!\mathit{omplexes}}
% \Pic is the operator that assigns to X its picard group, usage \Pic(X)
% \Picardstack_{X/B} denotes the Picard stack of X over B
% \Picardfunctor_{X/B} denotes the Picard functor of X over B
\def\Pic{\mathop{\mathrm{Pic}}\nolimits}
\def\Picardstack{\mathcal{P}\!\mathit{ic}}
\def\Picardfunctor{\mathrm{Pic}}
\def\Deformationcategory{\mathcal{D}\!\mathit{ef}}
\usepackage{cleveref}
\crefname{originalremark}{original remark}{original remarks}
\Crefname{originalremark}{Original remark}{Original remarks}
\makeatletter
\newlabel{formal-spaces-remark-permanence-variant-Noetherian}{{0GBP}{source}{}{sourcepointer-1}{}}
\newlabel{formal-spaces-remark-permanence-variant-Noetherian@cref}{{[originalremark][0][]0GBP}{[1][0][]source}{}{}{}}
\makeatother
\crefname{originalreference}{original reference}{original references}
\Crefname{originalreference}{Original reference}{Original references}
\makeatletter
\newlabel{lemma-equivalent-with-artin}{{0AJU}{source}{}{sourcepointer-2}{}}
\newlabel{lemma-equivalent-with-artin@cref}{{[originalreference][0][]0AJU}{[1][0][]source}{}{}{}}
\makeatother
\crefname{originalreference}{original reference}{original references}
\Crefname{originalreference}{Original reference}{Original references}
\makeatletter
\newlabel{lemma-exact-sequence-NL}{{00S2}{source}{}{sourcepointer-3}{}}
\newlabel{lemma-exact-sequence-NL@cref}{{[originalreference][0][]00S2}{[1][0][]source}{}{}{}}
\makeatother
\crefname{originalreference}{original reference}{original references}
\Crefname{originalreference}{Original reference}{Original references}
\makeatletter
\newlabel{lemma-rig-etale-continuous}{{0GCU}{source}{}{sourcepointer-4}{}}
\newlabel{lemma-rig-etale-continuous@cref}{{[originalreference][0][]0GCU}{[1][0][]source}{}{}{}}
\makeatother
\title{Rig-etale maps satisfy cancellation}
\author{The Stacks Project Authors}
\date{stacks-project:restricted.tex:local-snapshot-20261012; GNU FDL 1.2 or later}
\begin{document}
\maketitle
\noindent\textit{Editorial note.} All delivered mathematical text is the author's own, copied from the pinned Stacks Project chapter file restricted.tex; no statement or proof step was rewritten, reordered or removed.\par\medskip
\noindent\textit{Editorial note.} The counted proof is the author's complete argument: the proof fixes a common ideal of definition, writes the six-term exact sequence, transports the annihilation exponents from the two rig-etale maps through the sequence, and applies the rig-etale equivalence lemma to B -> C\par\medskip
\noindent\textit{Editorial note.} The label lemma-permanence-rig-etale-continuous names the counted statement, and the proof environment immediately following it in the source is delivered verbatim as its proof.\par\medskip
\noindent\textit{Editorial note.} The counted statement is carried by the author's own lemma environment in the chapter source.\par\medskip
\noindent\textit{Allowed standard prerequisites (editorial).}\begin{itemize}
\item the arrow property rig-etale for continuous homomorphisms of Noetherian adic rings (lemma-rig-etale-continuous of the same chapter) and the cancellation property of Formal Spaces, Remark formal-spaces-remark-permanence-variant-Noetherian
\item the exact sequence of completed naive cotangent complexes for A -> B -> C (lemma-exact-sequence-NL of the same chapter)
\item the rig-etale criterion in terms of annihilation of the cohomology of NL\textasciicircum{}wedge by powers of an ideal of definition (lemma-equivalent-with-artin of the same chapter)
\end{itemize}\medskip
\ifdefined\hypertarget\hypertarget{sourcepointer-1}{}\fi
\noindent\textit{Source reference.} Original remark 0GBP: sources/stacks-chow/tags.txt, source lines 20379--20379. Prerequisite pointer invoked at proof line(s) 4969; it names a standard fact of a sibling chapter that the counted argument uses, and every step of the counted proof built on it is retained verbatim in the delivered body.\par\medskip
\ifdefined\hypertarget\hypertarget{sourcepointer-2}{}\fi
\noindent\textit{Source reference.} Original reference 0AJU: sources/stacks-restricted/pinned-original-source.tex, source lines 1810--1810. Prerequisite pointer invoked at proof line(s) 5001; it names a standard fact of the same chapter that the counted argument uses, and every step of the counted proof built on it is retained verbatim in the delivered body.\par\medskip
\ifdefined\hypertarget\hypertarget{sourcepointer-3}{}\fi
\noindent\textit{Source reference.} Original reference 00S2: sources/stacks-restricted/pinned-original-source.tex, source lines 725--725. Prerequisite pointer invoked at proof line(s) 4981; it names a standard fact of the same chapter that the counted argument uses, and every step of the counted proof built on it is retained verbatim in the delivered body.\par\medskip
\ifdefined\hypertarget\hypertarget{sourcepointer-4}{}\fi
\noindent\textit{Source reference.} Original reference 0GCU: sources/stacks-restricted/pinned-original-source.tex, source lines 4749--4749. Prerequisite pointer invoked at proof line(s) 4974; it names a standard fact of the same chapter that the counted argument uses, and every step of the counted proof built on it is retained verbatim in the delivered body.\par\medskip
\begin{lemma}
\label{lemma-permanence-rig-etale-continuous}
The property $P(\varphi)=$``$\varphi$ is rig-\'etale''
on arrows of $\textit{WAdm}^{Noeth}$ has the cancellation property
as defined in Formal Spaces, Remark
\ref{formal-spaces-remark-permanence-variant-Noetherian}.
\end{lemma}

\begin{proof}
The statement makes sense by
Lemma \ref{lemma-rig-etale-continuous}.
To see that it is true assume we have maps
$A \to B$ and $B \to C$ in $\textit{WAdm}^{Noeth}$
with $A \to C$ and $A \to B$ rig-\'etale.
We have to show that $B \to C$ is rig-\'etale.
Then we can choose an ideal of definition $I \subset A$
such that the topology on $C$ and $B$ is the $I$-adic topology.
By Lemma \ref{lemma-exact-sequence-NL} we obtain an exact sequence
$$
\xymatrix{
C \otimes_B H^0(\NL_{B/A}^\wedge) \ar[r] &
H^0(\NL_{C/A}^\wedge) \ar[r] &
H^0(\NL_{C/B}^\wedge) \ar[r] & 0 \\
H^{-1}(\NL_{B/A}^\wedge \otimes_B C) \ar[r] &
H^{-1}(\NL_{C/A}^\wedge) \ar[r] &
H^{-1}(\NL_{C/B}^\wedge) \ar[llu]
}
$$
There exists a $c \geq 0$ such that for all $a \in I$ 
multiplication by $a^c$ is zero on $\NL_{B/A}^\wedge$ in $D(B)$ and
$\NL_{C/A}^\wedge$ in $D(C)$. Hence
$H^0(\NL_{B/A}^\wedge) \otimes_A C$,
$H^0(\NL_{C/A}^\wedge)$, and
$H^{-1}(\NL_{C/A}^\wedge)$
are annihilated by $a^c$. From the exact sequence
we obtain that multiplication by $a^{2c}$ is zero on
$H^0(\NL_{C/B}^\wedge)$ and $H^{-1}(\NL_{C/B}^\wedge)$.
It follows from Lemma \ref{lemma-equivalent-with-artin}
that $C$ is rig-\'etale over $(B, IB)$ as desired.
\end{proof}

\end{document}
